% --- Archon physics formalization source begin ---
% archon:physics
% archon:covers PhyXMiniProblems/problem_phyx_mini_0888.lean
% archon:source-report reports/phyx_mini/problem_phyx_mini_0888.source.json

\chapter{Physics problem phyx\_mini\_0888}
\label{ch:PhyXMiniProblems_problem_phyx_mini_0888}

\paragraph{Problem source.}
\#\# Physical scenario

analyze the RL circuit of figure.

\#\# Question

What is \textbackslash{}( V\_\textbackslash{}text\{R\} \textbackslash{}) in the limits \textbackslash{}( \textbackslash{}omega \textbackslash{}to 0 \textbackslash{})?

\#\# Answer choices

- A: A: \textbackslash{}( V\_\textbackslash{}text\{R\} \textbackslash{}to 2 \textbackslash{})

- B: B: \textbackslash{}( V\_\textbackslash{}text\{R\} \textbackslash{}to 0 \textbackslash{})

- C: C: \textbackslash{}( V\_\textbackslash{}text\{R\} \textbackslash{}to \textbackslash{}mathcal\{E\}\_0 \textbackslash{})

- D: D: \textbackslash{}( V\_\textbackslash{}text\{R\} \textbackslash{}to 3\textbackslash{})

\#\# Figure caption (auxiliary; use the image as primary evidence)

The image depicts a schematic of an R-L (resistor-inductor) AC circuit. 

- **Components and Labels:**
  - There's an AC voltage source on the left, represented by a circle with a sine wave symbol inside it.
  - The voltage source is labeled with \textbackslash{}( \textbackslash{}mathcal\{E\}\_0 \textbackslash{}cos \textbackslash{}omega t \textbackslash{}), indicating an alternating current voltage of amplitude \textbackslash{}( \textbackslash{}mathcal\{E\}\_0 \textbackslash{}), angular frequency \textbackslash{}( \textbackslash{}omega \textbackslash{}), and time \textbackslash{}( t \textbackslash{}).
  - A resistor is depicted at the top right, represented by a zigzag line, and labeled \textbackslash{}( R \textbackslash{}).
  - An inductor is shown below the resistor, represented by a coiled line, and labeled \textbackslash{}( L \textbackslash{}).

- **Connections:**
  - These components are connected in series: the voltage source connects to the resistor, which is then connected to the inductor, completing the circuit back to the voltage source. 

This image represents a basic R-L circuit driven by an AC source.

\#\# Dataset metadata

Electric Circuits; Physical Model Grounding Reasoning, Numerical Reasoning

\paragraph{Recorded answer/context.}
C: C: \textbackslash{}( V\_\textbackslash{}text\{R\} \textbackslash{}to \textbackslash{}mathcal\{E\}\_0 \textbackslash{})

\paragraph{Figure/image path.}
/root/proposal\_for\_physic/science-mango/phyx\_mini\_run/phyx\_data/test\_image/888.png

\paragraph{Formalization target.}
create a compiling Lean file with sorry bodies at `PhyXMiniProblems/problem\_phyx\_mini\_0888.lean`. The Lean declarations must preserve the physical quantities, dimensions or dimensional roles, figure labels, governing-law hypotheses, and final relation expressed by this problem.
Use Mathlib/Physlib names found through LeanExplore where available. If a physics API is missing, introduce faithful local abstractions rather than scalar placeholder aliases.

\begin{theorem}[Physics formalization target]
\label{thm:physics:phyx_mini_0888:target}
\lean{PhyXMiniProblems.ProblemPhyXMini0888.problem_phyx_mini_0888}
\uses{def:physics:phyx-mini-0888:phyxminiproblems-problemphyxmini0888-potentialamplitudeinvolts, def:physics:phyx-mini-0888:phyxminiproblems-problemphyxmini0888-seriesrlcircuitsetup, def:physics:phyx-mini-0888:phyxminiproblems-problemphyxmini0888-matchessuppliedseriesrlcircuitfigure, def:physics:phyx-mini-0888:phyxminiproblems-problemphyxmini0888-calibratesangularfrequencyparameter, def:physics:phyx-mini-0888:phyxminiproblems-problemphyxmini0888-hasphysicalseriesrlparameters, def:physics:phyx-mini-0888:phyxminiproblems-problemphyxmini0888-satisfiescosinesourcelaw, def:physics:phyx-mini-0888:phyxminiproblems-problemphyxmini0888-satisfiesseriesrlcircuitlaws}
The assigned autoformalize agent should translate this physics problem into Lean declarations in the covered file, with theorem and lemma proof bodies written as `by sorry`.
\end{theorem}
\begin{proof}
This is an autoformalization task, not a proof task. Produce statements that can later be proved by the physics prover without weakening the physical model.
\end{proof}

% --- Archon declaration topology begin ---
\section{Declaration topology}
\begin{definition}[electric Current Dimension]
  \label{def:physics:phyx-mini-0888:phyxminiproblems-problemphyxmini0888-electriccurrentdimension}
  \lean{PhyXMiniProblems.ProblemPhyXMini0888.electricCurrentDimension}
  Electric current has physical dimension charge per time
\end{definition}
\begin{proof}
  This is the declared model definition of electric Current Dimension.
\end{proof}
\begin{definition}[electric Potential Dimension]
  \label{def:physics:phyx-mini-0888:phyxminiproblems-problemphyxmini0888-electricpotentialdimension}
  \lean{PhyXMiniProblems.ProblemPhyXMini0888.electricPotentialDimension}
  Electric potential has physical dimension energy per charge
\end{definition}
\begin{proof}
  This is the declared model definition of electric Potential Dimension.
\end{proof}
\begin{definition}[electrical Resistance Dimension]
  \label{def:physics:phyx-mini-0888:phyxminiproblems-problemphyxmini0888-electricalresistancedimension}
  \lean{PhyXMiniProblems.ProblemPhyXMini0888.electricalResistanceDimension}
  \uses{def:physics:phyx-mini-0888:phyxminiproblems-problemphyxmini0888-electriccurrentdimension, def:physics:phyx-mini-0888:phyxminiproblems-problemphyxmini0888-electricpotentialdimension}
  Electrical resistance, and hence impedance magnitude, is voltage per current
\end{definition}
\begin{proof}
  This is the declared model definition of electrical Resistance Dimension.
\end{proof}
\begin{definition}[inductance Dimension]
  \label{def:physics:phyx-mini-0888:phyxminiproblems-problemphyxmini0888-inductancedimension}
  \lean{PhyXMiniProblems.ProblemPhyXMini0888.inductanceDimension}
  \uses{def:physics:phyx-mini-0888:phyxminiproblems-problemphyxmini0888-electriccurrentdimension, def:physics:phyx-mini-0888:phyxminiproblems-problemphyxmini0888-electricpotentialdimension}
  Inductance is voltage times time per current, as in V\_L = L dI/dt
\end{definition}
\begin{proof}
  This is the declared model definition of inductance Dimension.
\end{proof}
\begin{definition}[angular Frequency Dimension]
  \label{def:physics:phyx-mini-0888:phyxminiproblems-problemphyxmini0888-angularfrequencydimension}
  \lean{PhyXMiniProblems.ProblemPhyXMini0888.angularFrequencyDimension}
  Angular frequency has inverse-time dimension; radians are dimensionless
\end{definition}
\begin{proof}
  This is the declared model definition of angular Frequency Dimension.
\end{proof}
\begin{definition}[Time Quantity]
  \label{def:physics:phyx-mini-0888:phyxminiproblems-problemphyxmini0888-timequantity}
  \lean{PhyXMiniProblems.ProblemPhyXMini0888.TimeQuantity}
  A signed, unit-independent physical time
\end{definition}
\begin{proof}
  This is the declared model definition of Time Quantity.
\end{proof}
\begin{definition}[Angular Frequency Quantity]
  \label{def:physics:phyx-mini-0888:phyxminiproblems-problemphyxmini0888-angularfrequencyquantity}
  \lean{PhyXMiniProblems.ProblemPhyXMini0888.AngularFrequencyQuantity}
  \uses{def:physics:phyx-mini-0888:phyxminiproblems-problemphyxmini0888-angularfrequencydimension}
  A nonnegative, unit-independent angular frequency
\end{definition}
\begin{proof}
  This is the declared model definition of Angular Frequency Quantity.
\end{proof}
\begin{definition}[Electric Current Amplitude Quantity]
  \label{def:physics:phyx-mini-0888:phyxminiproblems-problemphyxmini0888-electriccurrentamplitudequantity}
  \lean{PhyXMiniProblems.ProblemPhyXMini0888.ElectricCurrentAmplitudeQuantity}
  \uses{def:physics:phyx-mini-0888:phyxminiproblems-problemphyxmini0888-electriccurrentdimension}
  A nonnegative, unit-independent electric-current amplitude
\end{definition}
\begin{proof}
  This is the declared model definition of Electric Current Amplitude Quantity.
\end{proof}
\begin{definition}[Electric Potential Quantity]
  \label{def:physics:phyx-mini-0888:phyxminiproblems-problemphyxmini0888-electricpotentialquantity}
  \lean{PhyXMiniProblems.ProblemPhyXMini0888.ElectricPotentialQuantity}
  \uses{def:physics:phyx-mini-0888:phyxminiproblems-problemphyxmini0888-electricpotentialdimension}
  A signed, unit-independent instantaneous electric potential
\end{definition}
\begin{proof}
  This is the declared model definition of Electric Potential Quantity.
\end{proof}
\begin{definition}[Electric Potential Amplitude Quantity]
  \label{def:physics:phyx-mini-0888:phyxminiproblems-problemphyxmini0888-electricpotentialamplitudequantity}
  \lean{PhyXMiniProblems.ProblemPhyXMini0888.ElectricPotentialAmplitudeQuantity}
  \uses{def:physics:phyx-mini-0888:phyxminiproblems-problemphyxmini0888-electricpotentialdimension}
  A nonnegative, unit-independent electric-potential amplitude
\end{definition}
\begin{proof}
  This is the declared model definition of Electric Potential Amplitude Quantity.
\end{proof}
\begin{definition}[Electrical Resistance Quantity]
  \label{def:physics:phyx-mini-0888:phyxminiproblems-problemphyxmini0888-electricalresistancequantity}
  \lean{PhyXMiniProblems.ProblemPhyXMini0888.ElectricalResistanceQuantity}
  \uses{def:physics:phyx-mini-0888:phyxminiproblems-problemphyxmini0888-electricalresistancedimension}
  A nonnegative, unit-independent electrical resistance
\end{definition}
\begin{proof}
  This is the declared model definition of Electrical Resistance Quantity.
\end{proof}
\begin{definition}[Inductance Quantity]
  \label{def:physics:phyx-mini-0888:phyxminiproblems-problemphyxmini0888-inductancequantity}
  \lean{PhyXMiniProblems.ProblemPhyXMini0888.InductanceQuantity}
  \uses{def:physics:phyx-mini-0888:phyxminiproblems-problemphyxmini0888-inductancedimension}
  A nonnegative, unit-independent inductance
\end{definition}
\begin{proof}
  This is the declared model definition of Inductance Quantity.
\end{proof}
\begin{definition}[Electrical Impedance Magnitude Quantity]
  \label{def:physics:phyx-mini-0888:phyxminiproblems-problemphyxmini0888-electricalimpedancemagnitudequantity}
  \lean{PhyXMiniProblems.ProblemPhyXMini0888.ElectricalImpedanceMagnitudeQuantity}
  \uses{def:physics:phyx-mini-0888:phyxminiproblems-problemphyxmini0888-electricalresistancedimension}
  A nonnegative, unit-independent impedance magnitude
\end{definition}
\begin{proof}
  This is the declared model definition of Electrical Impedance Magnitude Quantity.
\end{proof}
\begin{definition}[signed SIReadout]
  \label{def:physics:phyx-mini-0888:phyxminiproblems-problemphyxmini0888-signedsireadout}
  \lean{PhyXMiniProblems.ProblemPhyXMini0888.signedSIReadout}
  Coherent-SI readout of a signed dimensionful quantity
\end{definition}
\begin{proof}
  This is the declared model definition of signed SIReadout.
\end{proof}
\begin{definition}[nonnegative SIReadout]
  \label{def:physics:phyx-mini-0888:phyxminiproblems-problemphyxmini0888-nonnegativesireadout}
  \lean{PhyXMiniProblems.ProblemPhyXMini0888.nonnegativeSIReadout}
  Coherent-SI readout of a nonnegative dimensionful quantity
\end{definition}
\begin{proof}
  This is the declared model definition of nonnegative SIReadout.
\end{proof}
\begin{definition}[time In Seconds]
  \label{def:physics:phyx-mini-0888:phyxminiproblems-problemphyxmini0888-timeinseconds}
  \lean{PhyXMiniProblems.ProblemPhyXMini0888.timeInSeconds}
  \uses{def:physics:phyx-mini-0888:phyxminiproblems-problemphyxmini0888-timequantity, def:physics:phyx-mini-0888:phyxminiproblems-problemphyxmini0888-signedsireadout}
  Read time in seconds
\end{definition}
\begin{proof}
  This is the declared model definition of time In Seconds.
\end{proof}
\begin{definition}[angular Frequency In Radians Per Second]
  \label{def:physics:phyx-mini-0888:phyxminiproblems-problemphyxmini0888-angularfrequencyinradianspersecond}
  \lean{PhyXMiniProblems.ProblemPhyXMini0888.angularFrequencyInRadiansPerSecond}
  \uses{def:physics:phyx-mini-0888:phyxminiproblems-problemphyxmini0888-angularfrequencyquantity, def:physics:phyx-mini-0888:phyxminiproblems-problemphyxmini0888-nonnegativesireadout}
  Read angular frequency in radians per second
\end{definition}
\begin{proof}
  This is the declared model definition of angular Frequency In Radians Per Second.
\end{proof}
\begin{definition}[current Amplitude In Amperes]
  \label{def:physics:phyx-mini-0888:phyxminiproblems-problemphyxmini0888-currentamplitudeinamperes}
  \lean{PhyXMiniProblems.ProblemPhyXMini0888.currentAmplitudeInAmperes}
  \uses{def:physics:phyx-mini-0888:phyxminiproblems-problemphyxmini0888-electriccurrentamplitudequantity, def:physics:phyx-mini-0888:phyxminiproblems-problemphyxmini0888-nonnegativesireadout}
  Read an electric-current amplitude in amperes
\end{definition}
\begin{proof}
  This is the declared model definition of current Amplitude In Amperes.
\end{proof}
\begin{definition}[potential In Volts]
  \label{def:physics:phyx-mini-0888:phyxminiproblems-problemphyxmini0888-potentialinvolts}
  \lean{PhyXMiniProblems.ProblemPhyXMini0888.potentialInVolts}
  \uses{def:physics:phyx-mini-0888:phyxminiproblems-problemphyxmini0888-electricpotentialquantity, def:physics:phyx-mini-0888:phyxminiproblems-problemphyxmini0888-signedsireadout}
  Read a signed instantaneous electric potential in volts
\end{definition}
\begin{proof}
  This is the declared model definition of potential In Volts.
\end{proof}
\begin{definition}[potential Amplitude In Volts]
  \label{def:physics:phyx-mini-0888:phyxminiproblems-problemphyxmini0888-potentialamplitudeinvolts}
  \lean{PhyXMiniProblems.ProblemPhyXMini0888.potentialAmplitudeInVolts}
  \uses{def:physics:phyx-mini-0888:phyxminiproblems-problemphyxmini0888-electricpotentialamplitudequantity, def:physics:phyx-mini-0888:phyxminiproblems-problemphyxmini0888-nonnegativesireadout}
  Read an electric-potential amplitude in volts
\end{definition}
\begin{proof}
  This is the declared model definition of potential Amplitude In Volts.
\end{proof}
\begin{definition}[resistance In Ohms]
  \label{def:physics:phyx-mini-0888:phyxminiproblems-problemphyxmini0888-resistanceinohms}
  \lean{PhyXMiniProblems.ProblemPhyXMini0888.resistanceInOhms}
  \uses{def:physics:phyx-mini-0888:phyxminiproblems-problemphyxmini0888-electricalresistancequantity, def:physics:phyx-mini-0888:phyxminiproblems-problemphyxmini0888-nonnegativesireadout}
  Read an electrical resistance in ohms
\end{definition}
\begin{proof}
  This is the declared model definition of resistance In Ohms.
\end{proof}
\begin{definition}[inductance In Henries]
  \label{def:physics:phyx-mini-0888:phyxminiproblems-problemphyxmini0888-inductanceinhenries}
  \lean{PhyXMiniProblems.ProblemPhyXMini0888.inductanceInHenries}
  \uses{def:physics:phyx-mini-0888:phyxminiproblems-problemphyxmini0888-inductancequantity, def:physics:phyx-mini-0888:phyxminiproblems-problemphyxmini0888-nonnegativesireadout}
  Read an inductance in henries
\end{definition}
\begin{proof}
  This is the declared model definition of inductance In Henries.
\end{proof}
\begin{definition}[impedance Magnitude In Ohms]
  \label{def:physics:phyx-mini-0888:phyxminiproblems-problemphyxmini0888-impedancemagnitudeinohms}
  \lean{PhyXMiniProblems.ProblemPhyXMini0888.impedanceMagnitudeInOhms}
  \uses{def:physics:phyx-mini-0888:phyxminiproblems-problemphyxmini0888-electricalimpedancemagnitudequantity, def:physics:phyx-mini-0888:phyxminiproblems-problemphyxmini0888-nonnegativesireadout}
  Read an impedance magnitude in ohms
\end{definition}
\begin{proof}
  This is the declared model definition of impedance Magnitude In Ohms.
\end{proof}
\begin{definition}[Source Waveform Shape]
  \label{def:physics:phyx-mini-0888:phyxminiproblems-problemphyxmini0888-sourcewaveformshape}
  \lean{PhyXMiniProblems.ProblemPhyXMini0888.SourceWaveformShape}
  The waveform name printed beside the source symbol
\end{definition}
\begin{proof}
  This is the declared model definition of Source Waveform Shape.
\end{proof}
\begin{definition}[Source Formula Label]
  \label{def:physics:phyx-mini-0888:phyxminiproblems-problemphyxmini0888-sourceformulalabel}
  \lean{PhyXMiniProblems.ProblemPhyXMini0888.SourceFormulaLabel}
  \uses{def:physics:phyx-mini-0888:phyxminiproblems-problemphyxmini0888-sourcewaveformshape}
  Typed presentation data for the source label ℰ₀ cos(ω t)
\end{definition}
\begin{proof}
  This is the declared model definition of Source Formula Label.
\end{proof}
\begin{definition}[Sinusoidal Voltage Source]
  \label{def:physics:phyx-mini-0888:phyxminiproblems-problemphyxmini0888-sinusoidalvoltagesource}
  \lean{PhyXMiniProblems.ProblemPhyXMini0888.SinusoidalVoltageSource}
  \uses{def:physics:phyx-mini-0888:phyxminiproblems-problemphyxmini0888-timequantity, def:physics:phyx-mini-0888:phyxminiproblems-problemphyxmini0888-angularfrequencyquantity, def:physics:phyx-mini-0888:phyxminiproblems-problemphyxmini0888-electricpotentialquantity, def:physics:phyx-mini-0888:phyxminiproblems-problemphyxmini0888-electricpotentialamplitudequantity}
  The sinusoidal source and its independent voltage waveform
\end{definition}
\begin{proof}
  This is the declared model definition of Sinusoidal Voltage Source.
\end{proof}
\begin{definition}[Resistor]
  \label{def:physics:phyx-mini-0888:phyxminiproblems-problemphyxmini0888-resistor}
  \lean{PhyXMiniProblems.ProblemPhyXMini0888.Resistor}
  \uses{def:physics:phyx-mini-0888:phyxminiproblems-problemphyxmini0888-angularfrequencyquantity, def:physics:phyx-mini-0888:phyxminiproblems-problemphyxmini0888-electricpotentialamplitudequantity, def:physics:phyx-mini-0888:phyxminiproblems-problemphyxmini0888-electricalresistancequantity}
  The resistor marked R, including its independent voltage response
\end{definition}
\begin{proof}
  This is the declared model definition of Resistor.
\end{proof}
\begin{definition}[Inductor]
  \label{def:physics:phyx-mini-0888:phyxminiproblems-problemphyxmini0888-inductor}
  \lean{PhyXMiniProblems.ProblemPhyXMini0888.Inductor}
  \uses{def:physics:phyx-mini-0888:phyxminiproblems-problemphyxmini0888-inductancequantity}
  The inductor marked L
\end{definition}
\begin{proof}
  This is the declared model definition of Inductor.
\end{proof}
\begin{definition}[Circuit Component]
  \label{def:physics:phyx-mini-0888:phyxminiproblems-problemphyxmini0888-circuitcomponent}
  \lean{PhyXMiniProblems.ProblemPhyXMini0888.CircuitComponent}
  The three component identities encountered in one traversal of the loop
\end{definition}
\begin{proof}
  This is the declared model definition of Circuit Component.
\end{proof}
\begin{definition}[Series RLCircuit Figure]
  \label{def:physics:phyx-mini-0888:phyxminiproblems-problemphyxmini0888-seriesrlcircuitfigure}
  \lean{PhyXMiniProblems.ProblemPhyXMini0888.SeriesRLCircuitFigure}
  \uses{def:physics:phyx-mini-0888:phyxminiproblems-problemphyxmini0888-sourceformulalabel, def:physics:phyx-mini-0888:phyxminiproblems-problemphyxmini0888-circuitcomponent}
  Literal topology and annotations transcribed from image 888.png
\end{definition}
\begin{proof}
  This is the declared model definition of Series RLCircuit Figure.
\end{proof}
\begin{definition}[Series RLCircuit Setup]
  \label{def:physics:phyx-mini-0888:phyxminiproblems-problemphyxmini0888-seriesrlcircuitsetup}
  \lean{PhyXMiniProblems.ProblemPhyXMini0888.SeriesRLCircuitSetup}
  \uses{def:physics:phyx-mini-0888:phyxminiproblems-problemphyxmini0888-angularfrequencyquantity, def:physics:phyx-mini-0888:phyxminiproblems-problemphyxmini0888-electriccurrentamplitudequantity, def:physics:phyx-mini-0888:phyxminiproblems-problemphyxmini0888-electricalimpedancemagnitudequantity, def:physics:phyx-mini-0888:phyxminiproblems-problemphyxmini0888-sinusoidalvoltagesource, def:physics:phyx-mini-0888:phyxminiproblems-problemphyxmini0888-resistor, def:physics:phyx-mini-0888:phyxminiproblems-problemphyxmini0888-inductor, def:physics:phyx-mini-0888:phyxminiproblems-problemphyxmini0888-seriesrlcircuitfigure}
  The complete physical setup. The frequency-indexed current, impedance, and resistor voltage are observables, not definitions made from the answer. physicalFrequencyAt associates the nonnegative scalar used in the limit with a genuine dimensionful angular frequency
\end{definition}
\begin{proof}
  This is the declared model definition of Series RLCircuit Setup.
\end{proof}
\begin{definition}[Matches Supplied Series RLCircuit Figure]
  \label{def:physics:phyx-mini-0888:phyxminiproblems-problemphyxmini0888-matchessuppliedseriesrlcircuitfigure}
  \lean{PhyXMiniProblems.ProblemPhyXMini0888.MatchesSuppliedSeriesRLCircuitFigure}
  \uses{def:physics:phyx-mini-0888:phyxminiproblems-problemphyxmini0888-seriesrlcircuitsetup}
  The topology and printed annotations visible in the supplied bitmap
\end{definition}
\begin{proof}
  This is the declared model definition of Matches Supplied Series RLCircuit Figure.
\end{proof}
\begin{definition}[Calibrates Angular Frequency Parameter]
  \label{def:physics:phyx-mini-0888:phyxminiproblems-problemphyxmini0888-calibratesangularfrequencyparameter}
  \lean{PhyXMiniProblems.ProblemPhyXMini0888.CalibratesAngularFrequencyParameter}
  \uses{def:physics:phyx-mini-0888:phyxminiproblems-problemphyxmini0888-angularfrequencyinradianspersecond, def:physics:phyx-mini-0888:phyxminiproblems-problemphyxmini0888-seriesrlcircuitsetup}
  The scalar limit parameter is the coherent-SI readout of physical ω
\end{definition}
\begin{proof}
  This is the declared model definition of Calibrates Angular Frequency Parameter.
\end{proof}
\begin{definition}[Has Physical Series RLParameters]
  \label{def:physics:phyx-mini-0888:phyxminiproblems-problemphyxmini0888-hasphysicalseriesrlparameters}
  \lean{PhyXMiniProblems.ProblemPhyXMini0888.HasPhysicalSeriesRLParameters}
  \uses{def:physics:phyx-mini-0888:phyxminiproblems-problemphyxmini0888-potentialamplitudeinvolts, def:physics:phyx-mini-0888:phyxminiproblems-problemphyxmini0888-resistanceinohms, def:physics:phyx-mini-0888:phyxminiproblems-problemphyxmini0888-inductanceinhenries, def:physics:phyx-mini-0888:phyxminiproblems-problemphyxmini0888-seriesrlcircuitsetup}
  Positivity and nonnegativity conditions for a physical passive RL circuit
\end{definition}
\begin{proof}
  This is the declared model definition of Has Physical Series RLParameters.
\end{proof}
\begin{definition}[Satisfies Cosine Source Law]
  \label{def:physics:phyx-mini-0888:phyxminiproblems-problemphyxmini0888-satisfiescosinesourcelaw}
  \lean{PhyXMiniProblems.ProblemPhyXMini0888.SatisfiesCosineSourceLaw}
  \uses{def:physics:phyx-mini-0888:phyxminiproblems-problemphyxmini0888-timeinseconds, def:physics:phyx-mini-0888:phyxminiproblems-problemphyxmini0888-angularfrequencyinradianspersecond, def:physics:phyx-mini-0888:phyxminiproblems-problemphyxmini0888-potentialinvolts, def:physics:phyx-mini-0888:phyxminiproblems-problemphyxmini0888-potentialamplitudeinvolts, def:physics:phyx-mini-0888:phyxminiproblems-problemphyxmini0888-seriesrlcircuitsetup}
  The source formula shown in the figure: its instantaneous potential is ℰ₀ cos(ω t) for every physical frequency and time. This general waveform law does not constrain the requested resistor voltage
\end{definition}
\begin{proof}
  This is the declared model definition of Satisfies Cosine Source Law.
\end{proof}
\begin{definition}[Satisfies Series RLCircuit Laws]
  \label{def:physics:phyx-mini-0888:phyxminiproblems-problemphyxmini0888-satisfiesseriesrlcircuitlaws}
  \lean{PhyXMiniProblems.ProblemPhyXMini0888.SatisfiesSeriesRLCircuitLaws}
  \uses{def:physics:phyx-mini-0888:phyxminiproblems-problemphyxmini0888-angularfrequencyinradianspersecond, def:physics:phyx-mini-0888:phyxminiproblems-problemphyxmini0888-currentamplitudeinamperes, def:physics:phyx-mini-0888:phyxminiproblems-problemphyxmini0888-potentialamplitudeinvolts, def:physics:phyx-mini-0888:phyxminiproblems-problemphyxmini0888-resistanceinohms, def:physics:phyx-mini-0888:phyxminiproblems-problemphyxmini0888-inductanceinhenries, def:physics:phyx-mini-0888:phyxminiproblems-problemphyxmini0888-impedancemagnitudeinohms, def:physics:phyx-mini-0888:phyxminiproblems-problemphyxmini0888-seriesrlcircuitsetup}
  Standard steady-state amplitude laws for a series RL circuit. The impedance magnitude is sqrt (R² + (ωL)²), the source amplitude is I₀ |Z|, and the resistor amplitude is I₀ R. These are governing relations at every physical frequency, not the low-frequency limit asked for in the problem
\end{definition}
\begin{proof}
  This is the declared model definition of Satisfies Series RLCircuit Laws.
\end{proof}
\begin{definition}[Answer Choice]
  \label{def:physics:phyx-mini-0888:phyxminiproblems-problemphyxmini0888-answerchoice}
  \lean{PhyXMiniProblems.ProblemPhyXMini0888.AnswerChoice}
  Labels of the four displayed answer choices
\end{definition}
\begin{proof}
  This is the declared model definition of Answer Choice.
\end{proof}
\begin{definition}[displayed Limit In Volts]
  \label{def:physics:phyx-mini-0888:phyxminiproblems-problemphyxmini0888-answerchoice-displayedlimitinvolts}
  \lean{PhyXMiniProblems.ProblemPhyXMini0888.AnswerChoice.displayedLimitInVolts}
  \uses{def:physics:phyx-mini-0888:phyxminiproblems-problemphyxmini0888-answerchoice}
  Scalar voltage-limit expressions printed beside the choices. The unitless numbers in A and D are retained exactly as printed and interpreted as volt readouts; choice C repeats the source-amplitude symbol ℰ₀
\end{definition}
\begin{proof}
  This is the declared model definition of displayed Limit In Volts.
\end{proof}
\begin{definition}[recorded Dataset Answer]
  \label{def:physics:phyx-mini-0888:phyxminiproblems-problemphyxmini0888-recordeddatasetanswer}
  \lean{PhyXMiniProblems.ProblemPhyXMini0888.recordedDatasetAnswer}
  \uses{def:physics:phyx-mini-0888:phyxminiproblems-problemphyxmini0888-answerchoice}
  The answer label recorded in the supplied dataset metadata
\end{definition}
\begin{proof}
  This is the declared model definition of recorded Dataset Answer.
\end{proof}
\begin{lemma}[resistor Voltage Amplitude eq series RLFormula]
  \label{lem:physics:phyx-mini-0888:phyxminiproblems-problemphyxmini0888-resistorvoltageamplitude-eq-seriesrlformula}
  \lean{PhyXMiniProblems.ProblemPhyXMini0888.resistorVoltageAmplitude_eq_seriesRLFormula}
  \uses{def:physics:phyx-mini-0888:phyxminiproblems-problemphyxmini0888-potentialamplitudeinvolts, def:physics:phyx-mini-0888:phyxminiproblems-problemphyxmini0888-resistanceinohms, def:physics:phyx-mini-0888:phyxminiproblems-problemphyxmini0888-inductanceinhenries, def:physics:phyx-mini-0888:phyxminiproblems-problemphyxmini0888-seriesrlcircuitsetup, def:physics:phyx-mini-0888:phyxminiproblems-problemphyxmini0888-calibratesangularfrequencyparameter, def:physics:phyx-mini-0888:phyxminiproblems-problemphyxmini0888-hasphysicalseriesrlparameters, def:physics:phyx-mini-0888:phyxminiproblems-problemphyxmini0888-satisfiesseriesrlcircuitlaws}
  At a calibrated angular-frequency readout ω, the governing amplitude laws give the usual closed form for V\_R. This is an intermediate general-frequency relation, not the requested low-frequency answer
\end{lemma}
\begin{proof}
  The conclusion follows from the governing relations and figure readouts named in the statement of resistor Voltage Amplitude eq series RLFormula.
\end{proof}
% --- Archon declaration topology end ---
% --- Archon physics formalization source end ---
